\documentclass[11pt,fontset=fandol]{ctexart}
\usepackage[a4paper,margin=23mm]{geometry}
\usepackage{amsmath,amssymb,mathtools,bm,mathrsfs,esint}
\usepackage{graphicx,longtable,booktabs,array,calc}
\usepackage[normalem]{ulem}
\usepackage{xcolor,hyperref}
\hypersetup{hidelinks}
\setcounter{secnumdepth}{-1}
\newcounter{none}
\providecommand{\tightlist}{\setlength{\itemsep}{0pt}\setlength{\parskip}{0pt}}
\providecommand{\passthrough}[1]{#1}
\usepackage{amssymb,mathrsfs,esint,bm,newunicodechar}
\allowdisplaybreaks
\setlength{\emergencystretch}{3em}
\providecommand{\llangle}{\langle\!\langle}
\providecommand{\rrangle}{\rangle\!\rangle}
\newunicodechar{∞}{\ensuremath{\infty}}
\newunicodechar{∈}{\ensuremath{\in}}
\newunicodechar{∑}{\ensuremath{\sum}}
\newunicodechar{∫}{\ensuremath{\int}}
\newunicodechar{≤}{\ensuremath{\leq}}
\newunicodechar{≥}{\ensuremath{\geq}}
\newunicodechar{≠}{\ensuremath{\neq}}
\newunicodechar{→}{\ensuremath{\to}}
\newunicodechar{⇒}{\ensuremath{\Rightarrow}}
\newunicodechar{α}{\ensuremath{\alpha}}
\newunicodechar{β}{\ensuremath{\beta}}
\newunicodechar{γ}{\ensuremath{\gamma}}
\newunicodechar{λ}{\ensuremath{\lambda}}
\newunicodechar{φ}{\ensuremath{\varphi}}
\newunicodechar{ε}{\ensuremath{\varepsilon}}
\newunicodechar{σ}{\ensuremath{\sigma}}
\newunicodechar{θ}{\ensuremath{\theta}}
\newunicodechar{δ}{\ensuremath{\delta}}
\newunicodechar{η}{\ensuremath{\eta}}
\newunicodechar{π}{\ensuremath{\pi}}
\newunicodechar{ξ}{\ensuremath{\xi}}
\newunicodechar{Σ}{\ensuremath{\Sigma}}
\newunicodechar{Δ}{\ensuremath{\Delta}}
\newunicodechar{∂}{\ensuremath{\partial}}
\newunicodechar{∀}{\ensuremath{\forall}}
\newunicodechar{∃}{\ensuremath{\exists}}
\newunicodechar{√}{\ensuremath{\surd}}
\newunicodechar{−}{\ensuremath{-}}
\newunicodechar{μ}{\ensuremath{\mu}}
\newunicodechar{ρ}{\ensuremath{\rho}}
\newunicodechar{τ}{\ensuremath{\tau}}
\newunicodechar{ω}{\ensuremath{\omega}}
\newunicodechar{Ω}{\ensuremath{\Omega}}
\newunicodechar{χ}{\ensuremath{\chi}}
\newunicodechar{ψ}{\ensuremath{\psi}}
\newunicodechar{Γ}{\ensuremath{\Gamma}}
\newunicodechar{Φ}{\ensuremath{\Phi}}
\newunicodechar{Λ}{\ensuremath{\Lambda}}
\newunicodechar{ℝ}{\ensuremath{\mathbb{R}}}
\newunicodechar{ℤ}{\ensuremath{\mathbb{Z}}}
\newunicodechar{ℕ}{\ensuremath{\mathbb{N}}}
\newunicodechar{△}{\ensuremath{\triangle}}
\newunicodechar{⊗}{\ensuremath{\otimes}}
\begin{document}
\section{2022年杭州电子科技大学601数学分析考研真题}\label{ux5e74ux676dux5ddeux7535ux5b50ux79d1ux6280ux5927ux5b66601ux6570ux5b66ux5206ux6790ux8003ux7814ux771fux9898}

来源：\texttt{HDU/601数学分析/2022/2022年杭州电子科技大学601数学分析考研真题.pdf}。原
PDF 页序：1, 2。

\subsection{杭州电子科技大学2022年数学分析真题}\label{ux676dux5ddeux7535ux5b50ux79d1ux6280ux5927ux5b662022ux5e74ux6570ux5b66ux5206ux6790ux771fux9898}

\subsection{一.填空题(5×8=40分)}\label{ux4e00.ux586bux7a7aux98985840ux5206}

\begin{enumerate}
\def\labelenumi{\arabic{enumi}.}
\item
  \(\lim _ { n \to \infty } ( 1 + { \frac { 1 } { 2 } } + \cdots + { \frac { 1 } { n } } ) =\)
\item
  \(( x ^ { 2 } e ^ { - 2 x } ) ^ { ( n ) }\)
\end{enumerate}

3.设 \(0 < k < \frac { 1 } { e }\) ，方程 \(\ln x - k x = 0\) 在
\(( 0 , + \infty )\) 内根的个数为

\begin{enumerate}
\def\labelenumi{\arabic{enumi}.}
\setcounter{enumi}{3}
\tightlist
\item
  \(\int \frac { e ^ { x } - e ^ { - x } } { e ^ { x } + e ^ { - x } } d x =\)
\end{enumerate}

5.若正项级数 \(\sum _ { n = 1 } ^ { \infty } u _ { n }\) 收敛且所有项
\(u _ { n } \leq 3\) ，则级数
\(\sum _ { n = 1 } ^ { \infty } \frac { u _ { n } } { \pi - u _ { n } }\)
是否收敛？

6.给定方程 \(F ( x , y , z ) = 0\) ，则
\(\frac { \partial z } { \partial x } \frac { \partial x } { \partial y } \frac { \partial y } { \partial z } =\)

7.曲线
\(y = \ln ( 1 - x ^ { 2 } ) , x \in \left[ 0 , \frac { 1 } { 2 } \right]\)
的弧长为

8.设C是椭圆
\(\frac { x ^ { 2 } } { a ^ { 2 } } + \frac { y ^ { 2 } } { b ^ { 2 } } = 1 ( a > 0 , b > 0 )\)
，则
\(\oint _ { C } ( x ^ { 3 } + 2 y ) d x + ( 4 x - 3 y ^ { 2 } ) d y =\)

\subsection{二.计算}\label{ux4e8c.ux8ba1ux7b97}

1（5分）若 \(\lambda > \frac { 1 } { 2 }\) ，钞算函数

\[
f ( x , y ) = \left\{ \begin{array} { l l } { \displaystyle ( x ^ { 2 } + y ^ { 2 } ) ^ { \lambda } \sin \frac { 1 } { \sqrt { x ^ { 2 } + y ^ { 2 } } } , ( x , y ) \neq ( 0 , 0 ) } \\ { 0 , } & { ( x , y ) = ( 0 , 0 ) } \end{array} \right.
\]

在(0,0)的一阶偏导数 \(f _ { \mathrm { x } } ( 0 , 0 )\) 和
\(f _ { y } ( 0 , 0 )\) 以及
\(\lim _ { ( x , y ) \to ( 0 , 0 ) } f ( x , y )\)

2.求幂级数
\(\sum _ { n = 1 } ^ { \infty } { \frac { x ^ { n } } { \alpha ^ { n } + \beta ^ { n } } } ( \alpha > 0 , \beta > 0 )\)
的收敛域.

3.计算曲线积分
\(\int _ { L } { \frac { d x + d y } { \left| x \right| + \left| y \right| } }\)
，其中L为由A(1,0)出发经由B(0,1)再到C(-1,0)连成的

折线.

4.设V是由椭球面
\(\frac { x ^ { 2 } } { a ^ { 2 } } + \frac { y ^ { 2 } } { b ^ { 2 } } + \frac { z ^ { 2 } } { c ^ { 2 } } = 1 \left( a > 0 , b > 0 , c > 0 \right)\)
在第一卦限部分的切平面和三个坐标平面所围成的封闭区域的体积，求V取最小值时对应的切平面方程.三.(3x10=
30)

1.求出函数
\(f ( x ) = ( x ^ { 2 } + x - 6 ) \left| x ^ { 3 } - 4 x \right|\)
的全部不可导点.

2.给出函数列
\(f _ { n } ( x ) = { \frac { x ( \ln n ) ^ { \mu } } { n ^ { x } } } ( n = 2 , 3 , 4 , \cdots )\)
，试问：当μ为何值时， \(\{ f _ { n } ( x ) \}\) 在\([ 0 , + \infty )\)
上一致收敛.

3.比较
\(\int _ { - \infty } ^ { + \infty } { \frac { \sin { x } } { x } }\)
dx和 \(\int _ { - \pi } ^ { \pi } { \frac { \sin { x } } { x } } d x\)
的大小，并给出证明.

四.(5.×10=50分)

1.若c≠0，利用函数极限的定义证明
\(\lim _ { \stackrel { . . . } { x }  \epsilon } \frac { 1 } { x } = \frac { 1 } { e }\)

2.如果函数f(x)满足条件f'(0)=0且
\(f ^ { \prime } ( x ) + \left[ f ^ { \prime } ( x ) \right] ^ { 2 } = x\)
，证明(0,f(0))是函数f(x)的拐点.

3.设实常数a,b不同时为零，利用二重积分换元法证明

\[
\underset { x ^ { 2 } + y ^ { 2 } \leq 1 } { \iint } f ( a x \neq b y ) d x d y = 2 \int _ { - 1 } ^ { 1 } \sqrt { 1 - u ^ { 2 } } f ( \sqrt { a ^ { 2 } + b ^ { 2 } } u ) d u .
\]

4.设函数f（x）在 \(\left[ a , b \right] \left( b > a > 0 \right)\)
上连续，在(a,b)上可导且 \(f ( a ) \neq f ( b )\) .证明：存在
\(5 , 1 7 \in ( a , b )\) ，使得

\[
f ^ { \prime } ( \xi ) = ( a + b ) ( a ^ { 2 } + b ^ { 2 } ) \frac { f ^ { \prime } ( \eta ) } { 4 \eta ^ { 3 } } .
\]

5.若两个正数α和β使得不等式
\(( 1 + \frac { 1 } { n } ) ^ { n + \alpha } \leq e \leq ( 1 + \frac { 1 } { n } ) ^ { n + \beta }\)
对所有的正整数n成立，证明α的最大值为
\(\alpha = \frac { 1 } { \ln 2 } - 1\) ，β的最小值为
\(\beta = \frac { 1 } { 2 }\)
\end{document}
